\chapter{Suites arithmétiques et géométriques}

\section{Suites arithmétiques}

\subsection{Définition}

\begin{ex}
	Considérons la suite $(u_{n})$ où l'on passe d'un terme au suivant en ajoutant $5$.

	Si le premier terme est égal à $3$, les termes suivants sont :

	\begin{itemize}
		\item $u_{\pa{0}} = 3$
		\item $u_{\pa{1}} = 8$
		\item $u_{\pa{2}} = 13$
		\item $u_{\pa{3}} = 18$
	\end{itemize}

	Une telle suite est appelée une suite \textbf{arithmétique} de raison $5$ et de premier terme $3$.

	La suite est définie par :
	\[
		\pa{u_n}:\left\{ \begin{array}{l}u_{\pa{0}} = 3\\u_{\pa{n + 1}} = u_{\pa{n}} + 5\end{array}\right.
	\]
\end{ex}

\begin{ex}
	Soit la suite \textbf{arithmétique} $(v_{n})$ de premier terme $5$ et de raison $-2$.

	Les premiers termes successifs sont :

	\begin{itemize}
		\item $v_{\pa{0}}=5$
		\item $v_{\pa{1}}=5-2 = 3$
		\item $v_{\pa{2}}=3-2 = 1$
		\item $v_{\pa{3}}=1-2 = -1$
	\end{itemize}

	La suite est donc définie par :
	\[
		(v_n):\left\{ \begin{array}{l}v_{\pa{0}} = 5 \\ v_{\pa{n + 1}} = v_{\pa{n}}-2\end{array} \right.
	\]
\end{ex}

\begin{definition}
	Une suite $(u_{n})$ est une \textbf{suite arithmétique} s'il existe un nombre $r$ tel que pour tout entier $n$, on a :

	\[
		u_{\pa{n + 1}} = u_{\pa{n}} + r
	\]

	Le nombre $r$ est appelé \textbf{raison} de la suite.

	\begin{center}
		\begin{tikzpicture}[xscale=1.3, yscale=0.9, >=Triangle]
			\draw[thick] (0,0.2) rectangle (10,1);
			\foreach \x in {1, 2, 3, 4, 5, 6, 7, 8, 9} { \draw[thick] (\x,0.2 ) -- (\x,1); }
			\node at (0.5, 0.55) {$u_{\pa{0}}$};
			\node at (1.5, 0.55) {$u_{\pa{1}}$};
			\node at (2.5, 0.55) {$u_{\pa{2}}$};
			\node at (3.5, 0.55) {$u_{\pa{3}}$};
			\node at (4.5, 0.55) {$\dots$};
			\node at (5.5, 0.55) {$\dots$};
			\node at (6.5, 0.55) {$u_{\pa{n-1}}$};
			\node at (7.5, 0.55) {$u_{\pa{n}}$};
			\node at (8.5, 0.55) {$u_{\pa{n+1}}$};
			\node at (9.5, 0.55) {$\dots$};
			\foreach \x in {0.5, 1.5, 2.5, 3.5, 5.5, 6.5, 7.5, 8.5} { \draw[->, very thick, blue!40!black] (\x, 1.2) to[out=70, in=140] node[midway, above=1pt] {$+r$} (\x+0.8, 1.2); }
		\end{tikzpicture}
	\end{center}
\end{definition}

\begin{ex}
	\[
		(u_n) : \begin{cases} u_{\pa{0}} = 10 \\ u_{\pa{n+1}} = u_{\pa{n}} \color{red}{+3} \end{cases}
		\qquad\qquad
		(v_n) : \begin{cases} v_{\pa{0}} = -15 \\ v_{\pa{n+1}} = v_{\pa{n}} \color{red}{-6} \end{cases}
		\qquad\qquad
		(w_n) : \begin{cases} w_{\pa{0}} = 300 \\ w_{\pa{n+1}} = w_{\pa{n}} \color{red}{+2{,}1} \end{cases}
	\]
\end{ex}

\subsection{Variations}

\begin{propriete}
	Soit $(u_{n})$ une suite \textbf{arithmétique} de raison $r$.

	\begin{itemize}
		\item Si $r>0$ alors la suite $(u_{n})$ est \textbf{croissante}.
		\item Si $r=0$ alors la suite $(u_{n})$ est \textbf{constante}.
		\item Si $r < 0$ alors la suite $(u_{n})$ est \textbf{décroissante}.
	\end{itemize}
\end{propriete}

\begin{demonstration}
	Soit $(u_{n})$ une suite \textbf{arithmétique} de raison $r$ :

	On a :
	\[
		u_{\pa{n + 1}}-u_{\pa{n}} = \pa{u_{\pa{n}} + r}- u_{\pa{n}} = r
	\]

	\begin{itemize}
		\item Si $r > 0$ alors $u_{\pa{n + 1}} - u_{\pa{n}} > 0$ et la suite $(u_{n})$ est \textbf{croissante}.
		\item Si $r < 0$ alors $u_{\pa{n + 1}} - u_{\pa{n}} < 0$ et la suite $(u_{n})$ est \textbf{décroissante}.
	\end{itemize}
\end{demonstration}

\begin{methode}[Déterminer le sens de variation d'une suite arithmétique]
	Étudions les variations des suites arithmétiques $(u_{n})$ et $(v_{n})$ définies par :

	\[
		u_{\pa{n}} = 3 + 5n \quad\text{et}\quad (v_n):\left\{ \begin{array}{l}v_{\pa{0}} = -3 \\ v_{\pa{n + 1}} = v_{\pa{n}}-4\end{array} \right.
	\]

	\begin{correction}
		\nline
		\begin{itemize}
			\item Soit $u_{\pa{n}} = 3 + 5n$.

			      En calculant quelques termes de $(u_n)$, on a :

			      \begin{center}
				      \rnc{\arraystretch}{1.2}
				      \begin{tabular}{|c|c|c|c|c|c|c|c|}
					      \hline
					      \ccg $n$   & $0$ & $1$ & $2$  & $3$  & $4$  & $5$  & $\dots$ \\ \hline
					      \ccg $u_n$ & $3$ & $8$ & $13$ & $18$ & $23$ & $28$ & $\dots$ \\ \hline
				      \end{tabular}
			      \end{center}

			      Donc la suite $(u_{n})$ est une suite \textbf{arithmétique croissante} de raison $r=5$ (positive).

			\item Soit la suite $(v_n)$ définie par $\left\{ \begin{array}{l}v_{\pa{0}} = -3 \\ v_{\pa{n + 1}} = v_{\pa{n}}\color{red}{-4}\end{array} \right.$.

			      On passe d'un terme au suivant en ajoutant $-4$.

			      Donc $(v_{n})$ est une suite \textbf{arithmétique décroissante} de raison $r=-4$ (négative).
		\end{itemize}
	\end{correction}
\end{methode}

\subsection{Représentation graphique}

\begin{propriete}
	Les points de la représentation graphique d'une suite arithmétique sont \textbf{alignés}.
\end{propriete}

\begin{ex}
	Soit $\pa{u_n}$ une suite arithmétique de raison $-0{,}5$ et de premier terme $4$.

	Chaque terme s'obtient en ajoutant la raison au terme précédent ($u_{\pa{n+1}} = u_{\pa{n}} - 0{,}5$) :

	\begin{itemize}
		\item $u_{\pa{0}} = 4$
		\item $u_{\pa{1}} = u_{\pa{0}} - 0{,}5 = 4 - 0{,}5 = 3{,}5$
		\item $u_{\pa{2}} = u_{\pa{1}} - 0{,}5 = 3{,}5 - 0{,}5 = 3$
		\item $u_{\pa{3}} = u_{\pa{2}} - 0{,}5 = 3 - 0{,}5 = 2{,}5$
	\end{itemize}

	On a donc :

	\begin{center}
		\rnc{\arraystretch}{1.4}
		\begin{tabular}{|c|c|c|c|c|c|c|c|}\hline
			\ccg $n$   & $0$ & $1$     & $2$ & $3$     & $4$ & $\dots$ & $12$ \\ \hline
			\ccg $u_n$ & $4$ & $3{,}5$ & $3$ & $2{,}5$ & $2$ & $\dots$ & $-2$ \\ \hline
		\end{tabular}
	\end{center}

	On a représenté ci-dessous $\pa{u_n}$ :

	\begin{center}
		\begin{tikzpicture}[scale=0.85, >=Stealth]
			\tikzset{ pt/.style={draw=blue!80!black, fill=blue!60, circle, inner sep=1.4pt} }
			\draw[gray!30, very thin, step=1] (-0.5,-2.5) grid (12.5,4.5);
			\draw[->, thick] (-0.5,0) -- (12.5,0) node[right] {$n$};
			\draw[->, thick] (0,-2.5) -- (0,4.5) node[above] {$u_n$};
			\foreach \x in {1,...,12} \draw (\x,0.1) -- (\x,-0.1) node[below, scale=0.8] {$\x$};
			\foreach \y in {-2,-1,1,2,3,4} \draw (0.1,\y) -- (-0.1,\y) node[left, scale=0.8] {$\y$};
			\node[below left, scale=0.8] at (0,0) {$0$};
			\foreach \n in {0, 1,..., 12} { \node[pt] at (\n, 4-0.5*\n) {}; }
		\end{tikzpicture}
	\end{center}
\end{ex}

\subsection{Expression du terme général}

\begin{propriete}
	Soit $(u_{n})$ une suite \textbf{arithmétique} de raison $r$. On a :

	\[
		u_{\pa{n}}=u_{\pa{0}}+n\times r
	\]

	\begin{center}
		\begin{tikzpicture}[xscale=1.3, yscale=0.9, >=Triangle]
			\draw[thick] (0,0.2) rectangle (10,1);
			\foreach \x in {1, 2, 3, 4, 5, 6, 7, 8, 9} { \draw[thick] (\x,0.2) -- (\x,1); }
			\node at (0.5, 0.55) {$u_{\pa{0}}$};
			\node at (1.5, 0.55) {$u_{\pa{1}}$};
			\node at (2.5, 0.55) {$u_{\pa{2}}$};
			\node at (3.5, 0.55) {$u_{\pa{3}}$};
			\node at (4.5, 0.55) {$\dots$};
			\node at (5.5, 0.55) {$\dots$};
			\node at (6.5, 0.55) {$u_{\pa{n-1}}$};
			\node at (7.5, 0.55) {$u_{\pa{n}}$};
			\node at (8.5, 0.55) {$u_{\pa{n+1}}$};
			\node at (9.5, 0.55) {$\dots$};
			\foreach \x in {0.5, 1.5, 2.5, 3.5, 5.5, 6.5, 7.5, 8.5} { \draw[->, very thick, blue!40!black] (\x, 1.2) to[out=70, in=140] node[midway, above=1pt] {$+r$} (\x+0.8, 1.2); }
			\draw[->, very thick, blue!40!black] (0.4, 0.0) .. controls (0.4,-1.0) and (2.5,-1.0) .. node[midway, fill=white, inner sep=1pt, scale=0.9] {$+(2 \times r)$} (2.5, 0.0);
			\draw[->, very thick, blue!40!black] (0.4, -0.1) .. controls (0.4,-2.0) and (3.5,-2.0) .. node[pos=0.6, fill=white, inner sep=1pt, scale=0.9] {$+(3 \times r)$} (3.5, 0.0);
			\draw[->, very thick, blue!40!black] (0.4, -0.2) .. controls (0.4,-3.2) and (7.5,-3.2) .. node[midway, fill=white, inner sep=2pt] {$+(n \times r)$} (7.5, 0.0);
		\end{tikzpicture}
	\end{center}
\end{propriete}

\begin{remarque}
	On a aussi : $u_{\pa{n}}=u_{\pa{\textcolor{blue}{1}}}+(n-\textcolor{blue}{1})\times r$
\end{remarque}

\begin{ex}
	Considérons une suite définie par : $u_{\pa{n+1}}=u_{\pa{n}}-3$ et $u_{\pa{0}}=10$.

	$(u_n)$ est une \textbf{suite arithmétique} de raison $r=-3$.

	\begin{center}
		\rnc{\arraystretch}{1.4}
		\begin{tabular}{|c|c|c|c|c|c|c|c|c|c|c}\hline
			\ccg $n$   & $0$  & $1$ & $2$ & $3$ & $4$  & $5$  & $6$  & $\dots$ & \ccg $\textcolor{red}{17}$ & $\dots$ \\ \hline
			\ccg $u_n$ & $10$ & $7$ & $4$ & $1$ & $-2$ & $-5$ & $-8$ & $\dots$ & \ccg $-41$                 & $\dots$ \\ \hline
		\end{tabular}
	\end{center}

	On a :
	\[
		u_{\pa{n}}=u_{\pa{0}}+n\times r \quad\Lrarr\quad u_{\pa{n}}=10-3n
	\]

	Avec cette expression, on a :
	\[
		u_{\pa{\textcolor{red}{17}}}=10-3\times \textcolor{red}{17} = 10-51 = -41
	\]
\end{ex}

\section{Suites géométriques}

\subsection{Définition}

\begin{ex}
	Considérons la suite $(u_{n})$ où l'on passe d'un terme au suivant en multipliant par $2$.

	Si le premier terme est égal à $5$, les termes suivants sont :

	\begin{itemize}
		\item $u_{\pa{0}} = 5$
		\item $u_{\pa{1}} = 10$
		\item $u_{\pa{2}} = 20$
		\item $u_{\pa{3}} = 40$
	\end{itemize}

	Une telle suite est appelée une suite \textbf{géométrique} de raison $2$ et de premier terme $5$.

	La suite $(u_n)$ est donc définie par : $\qquad\pa{u_n}:\begin{cases}u_{\pa{0}} = 5 \\u_{\pa{n + 1}} = 2u_{\pa{n}}\end{cases}$ \end{ex}

\begin{ex}
	Soit la suite géométrique $(v_{n})$ de premier terme $4$ et de raison $0{,}1$.

	Les premiers termes successifs sont :

	\begin{itemize}
		\item $v_{\pa{0}}=4$
		\item $v_{\pa{1}}=0{,}1 \times 4 = 0{,}4$
		\item $v_{\pa{2}}=0{,}1 \times 0{,}4 = 0{,}04$
		\item $v_{\pa{3}}=0{,}1 \times 0{,}04 = 0{,}004$
	\end{itemize}

	La suite $(v_n)$ est donc définie par : $\qquad\pa{v_n}:\begin{cases}v_{\pa{0}} = 4 \\ v_{\pa{n + 1}} = 0{,}1 \times v_{\pa{n}}\end{cases}$
\end{ex}

\begin{definition}
	Une suite $(u_{n})$ est une \textbf{suite géométrique} s'il existe un nombre $q$, \textbf{strictement positif}, tel que, pour tout entier $n$, on a :

	\[
		u_{\pa{n + 1}} = u_{\pa{n}}\times q
	\]

	Le nombre $q$ est appelé \textbf{raison} de la suite.

	\begin{center}
		\begin{tikzpicture}[xscale=1.3, yscale=0.9, >=Triangle]
			\draw[thick] (0,0.2) rectangle (10,1);
			\foreach \x in {1, 2, 3, 4, 5, 6, 7, 8, 9} { \draw[thick] (\x,0.2) -- (\x,1); }
			\node at (0.5, 0.55) {$u_{\pa{0}}$};
			\node at (1.5, 0.55) {$u_{\pa{1}}$};
			\node at (2.5, 0.55) {$u_{\pa{2}}$};
			\node at (3.5, 0.55) {$u_{\pa{3}}$};
			\node at (4.5, 0.55) {$\dots$};
			\node at (5.5, 0.55) {$\dots$};
			\node at (6.5, 0.55) {$u_{\pa{n-1}}$};
			\node at (7.5, 0.55) {$u_{\pa{n}}$};
			\node at (8.5, 0.55) {$u_{\pa{n+1}}$};
			\node at (9.5, 0.55) {$\dots$};
			\foreach \x in {0.5, 1.5, 2.5, 3.5, 5.5, 6.5, 7.5, 8.5} { \draw[->, very thick, blue!40!black] (\x, 1.2) to[out=70, in=140] node[midway, above=1pt] {$\times q$} (\x+0.8, 1.2); }
		\end{tikzpicture}
	\end{center}
\end{definition}

\begin{ex}
	\[
		(u_n) : \begin{cases} u_{\pa{0}} = 2 \\ u_{\pa{n+1}} = {\color{red}{7}}\times u_{\pa{n}} \end{cases}
		\qquad
		(v_n) : \begin{cases} v_{\pa{0}} = 1\,500 \\ v_{\pa{n+1}} = {\color{red}{\dfrac{2}{5}}}\times v_{\pa{n}} \end{cases}
		\qquad
		(w_n) : \begin{cases} w_{\pa{0}} = 300 \\ w_{\pa{n+1}} = {\color{red}{0{,}4}}\times w_{\pa{n}} \end{cases}
	\]
\end{ex}

\begin{ex}
	On place un capital de $500\eur{}$ sur un compte dont les intérêts annuels s'élèvent à $4\,\%$.

	Chaque année, le capital est donc multiplié par $1{,}04$.

	Ce capital suit une \textbf{progression géométrique} de raison $1{,}04$.

	On a ainsi :

	\begin{itemize}
		\item $u_{\pa{1}} = 1{,}04 \times 500 = 520$
		\item $u_{\pa{2}} = 1{,}04 \times 520 = 540{,}80$
		\item $u_{\pa{3}} = 1{,}04 \times 540{,}80 = 562{,}432$
	\end{itemize}

	De manière générale : $u_{\pa{n + 1}} = 1{,}04 \times u_{\pa{n}}$ avec $u_{\pa{0}} = 500$.
\end{ex}

\subsection{Variations}

\begin{propriete}
	Soit $(u_{n})$ une suite \textbf{géométrique} de raison $q$ et de premier terme $u_{\pa{0}}$ \textbf{strictement positif}.

	\begin{itemize}
		\item Si $q>1$ alors la suite $(u_{n})$ est \textbf{croissante}.
		\item Si $q=1$ alors la suite $(u_{n})$ est \textbf{constante}.
		\item Si $0 < q < 1$ alors la suite $(u_{n})$ est \textbf{décroissante}.
	\end{itemize}
\end{propriete}

\begin{methode}[Déterminer le sens de variation d'une suite géométrique]
	Déterminons le sens de variation des suites géométriques $(u_{n})$ et $(v_{n})$ définies par :

	\[
		u_{\pa{n}} = 4 \times 2^{n} \quad\text{et}\quad (v_n):\left\{\begin{array}{l}v_{\pa{0}} = 2 \\v_{\pa{n + 1}} = \dfrac{1}{2}v_{\pa{n}}\end{array} \right.
	\]

	\begin{correction}
		\nline
		\begin{itemize}
			\item Soit $u_{\pa{n}} = 4 \times 2^{n}$. En calculant quelques termes de $(u_n)$, on a :

			      \begin{center}
				      \rnc{\arraystretch}{1.2}
				      \begin{tabular}{|c|c|c|c|c|c|c|c|}
					      \hline
					      \ccg $n$   & $0$ & $1$ & $2$  & $3$  & $4$  & $5$   & $\dots$ \\ \hline
					      \ccg $u_n$ & $4$ & $8$ & $16$ & $32$ & $64$ & $128$ & $\dots$ \\ \hline
				      \end{tabular}
			      \end{center}

			      La suite $(u_{n})$ est géométrique et croissante car $q = 2 > 1$.

			\item Soit la suite $(v_{n})$ définie par $\left\{\begin{array}{l}v_{\pa{0}} = 2 \\ v_{\pa{n + 1}} = {\color{red}\dfrac{1}{2}}v_{\pa{n}}\end{array}\right.$.

			      Donc $(v_n)$ est \textbf{décroissante} car $q = {\color{red}\dfrac{1}{2}}$ est compris entre $0$ et $1$.

			      Calcul de quelques termes :

			      \begin{center}
				      \rnc{\arraystretch}{1.2}
				      \begin{tabular}{|c|c|c|c|c|c|c|}
					      \hline
					      \ccg $n$   & $0$ & $1$ & $2$     & $3$      & $4$       & $\dots$ \\ \hline
					      \ccg $v_n$ & $2$ & $1$ & $0{,}5$ & $0{,}25$ & $0{,}125$ & $\dots$ \\ \hline
				      \end{tabular}
			      \end{center}
		\end{itemize}
	\end{correction}
\end{methode}

\subsection{Représentation graphique}

\begin{propriete}
	Les points de la représentation graphique d'une \textbf{suite géométrique} ne sont pas \textbf{alignés}.

	\begin{center}
		\begin{tikzpicture}
			\tikzset{
				axe/.style={->, >=Stealth, thick, black!70},
				pt/.style={draw=blue!80!black, fill=blue!60, circle, inner sep=1.4pt},
				cadre/.style={draw=gray!50, rounded corners=2pt, fill=white}
			}
			\begin{scope}[shift={(0,1.5)}]
				\draw[cadre] (-0.5,-0.3) rectangle (4.5,2.8);
				\draw[axe] (0,0) -- (4,0) node[right, scale=0.8] {$n$};
				\draw[axe] (0,0) -- (0,2.4) node[above, scale=0.8] {$u_n$};
				\node[scale=0.95, anchor=west] at (0.3,2.1) {\textbf{$u_0 > 0$ et $0 < q < 1$}};
				\foreach \n in {0, 1, 2, 3, 4, 5} { \node[pt] at (\n*0.7, 2*0.5^\n) {}; }
			\end{scope}
			\begin{scope}[shift={(5.5,1.5)}]
				\draw[cadre] (-0.5,-0.3) rectangle (4.5,2.8);
				\draw[axe] (0,0) -- (4,0) node[right, scale=0.8] {$n$};
				\draw[axe] (0,0) -- (0,2.4) node[above, scale=0.8] {$u_n$};
				\node[scale=0.95, anchor=west] at (0.3,2.1) {\textbf{$u_0 > 0$ et $q > 1$}};
				\foreach \n in {0, 1, 2, 3, 4, 5} { \node[pt] at (\n*0.7, 0.1*1.8^\n) {}; }
			\end{scope}
			\begin{scope}[shift={(0,0)}]
				\draw[cadre] (-0.5,-2.6) rectangle (4.5,0.9);
				\draw[axe] (0,0) -- (4,0) node[right, scale=0.8] {$n$};
				\draw[axe] (0,-2.4) -- (0,0.5) node[above, scale=0.8] {$u_n$};
				\node[scale=0.95, anchor=west] at (0.3,-2.1) {\textbf{$u_0 < 0$ et $0 < q < 1$}};
				\foreach \n in {0, 1, 2, 3, 4, 5} { \node[pt] at (\n*0.7, -2.2*0.55^\n) {}; }
			\end{scope}
			\begin{scope}[shift={(5.5,0)}]
				\draw[cadre] (-0.5,-2.6) rectangle (4.5,0.9);
				\draw[axe] (0,0) -- (4,0) node[right, scale=0.8] {$n$};
				\draw[axe] (0,-2.4) -- (0,0.5) node[above, scale=0.8] {$u_n$};
				\node[scale=0.95, anchor=west] at (0.3,-2.1) {\textbf{$u_0 < 0$ et $q > 1$}};
				\foreach \n/\v in {0/-0.15, 1/-0.26, 2/-0.45, 3/-0.78, 4/-1.36, 5/-2.37} {
						\node[pt] at (\n*0.7, 0.9*\v) {};
					}
			\end{scope}
		\end{tikzpicture}
	\end{center}
\end{propriete}

\subsection{Expression du terme général}

\begin{propriete}
	Soit $(u_{n})$ une suite \textbf{géométrique} de raison $q$. On a :

	\[
		u_{\pa{n}}=u_{\pa{0}}\times q^n
	\]

	\begin{center}
		\begin{tikzpicture}[xscale=1.3, yscale=0.9, >=Triangle]
			\draw[thick] (0,0.2) rectangle (10,1);
			\foreach \x in {1, 2, 3, 4, 5, 6, 7, 8, 9} { \draw[thick] (\x,0.2) -- (\x,1); }
			\node at (0.5, 0.55) {$u_{\pa{0}}$};
			\node at (1.5, 0.55) {$u_{\pa{1}}$};
			\node at (2.5, 0.55) {$u_{\pa{2}}$};
			\node at (3.5, 0.55) {$u_{\pa{3}}$};
			\node at (4.5, 0.55) {$\dots$};
			\node at (5.5, 0.55) {$\dots$};
			\node at (6.5, 0.55) {$u_{\pa{n-1}}$};
			\node at (7.5, 0.55) {$u_{\pa{n}}$};
			\node at (8.5, 0.55) {$u_{\pa{n+1}}$};
			\node at (9.5, 0.55) {$\dots$};
			\foreach \x in {0.5, 1.5, 2.5, 3.5, 5.5, 6.5, 7.5, 8.5} { \draw[->, very thick, blue!40!black] (\x, 1.2) to[out=70, in=140] node[midway, above=1pt] {$\times q$} (\x+0.8, 1.2); }
			\draw[->, very thick, blue!40!black] (0.4, 0.0) .. controls (0.4,-1.0) and (2.5,-1.0) .. node[midway, fill=white, inner sep=1pt, scale=0.9] {$\times q^2$} (2.5, 0.0);
			\draw[->, very thick, blue!40!black] (0.4, -0.1) .. controls (0.4,-2.0) and (3.5,-2.0) .. node[pos=0.6, fill=white, inner sep=1pt, scale=0.9] {$\times q^3$} (3.5, 0.0);
			\draw[->, very thick, blue!40!black] (0.4, -0.2) .. controls (0.4,-3.2) and (7.5,-3.2) .. node[midway, fill=white, inner sep=2pt] {$\times q^n$} (7.5, 0.0);
		\end{tikzpicture}
	\end{center}
\end{propriete}

\begin{remarque}
	On a aussi : $u_{\pa{n}}=u_{\pa{\textcolor{blue}{1}}}\times q^{(n-\textcolor{blue}{1})}$
\end{remarque}

\begin{ex}
	Considérons une suite définie par : $u_{\pa{n+1}}=u_{\pa{n}}\times 0{,}25$ et $u_{\pa{0}}=100$.

	$(u_n)$ est une \textbf{suite géométrique} de raison $q=0{,}25$.

	\begin{center}
		\rnc{\arraystretch}{1.4}
		\begin{tabular}{|c|c|c|c|c|c|c|c|c}\hline
			\ccg $n$   & $0$   & $1$  & $2$      & $3$        & $4$          & $\dots$ & \ccg $\textcolor{red}{7}$ & $\dots$ \\ \hline
			\ccg $u_n$ & $100$ & $25$ & $6{,}25$ & $1{,}5625$ & $0{,}390625$ & $\dots$ & \ccg $0{,}0061\dots$      & $\dots$ \\ \hline
		\end{tabular}
	\end{center}

	On a :
	\[
		u_{\pa{n}}=u_{\pa{0}}\times q^{n} \quad\Lrarr\quad u_{\pa{n}}=100\times 0{,}25^{n}
	\]

	Avec cette expression :
	\[
		u_{\pa{\textcolor{red}{7}}} = 100\times 0{,}25^{\textcolor{red}{7}} \approx 0{,}0061
	\]
\end{ex}

\begin{ex}
	Un capital de $2\,500\eur{}$ est placé à un taux d'intérêt annuel de $7{,}5\,\%$.

	On note $u_{\pa{n}}$ le capital au bout de $n$ années.

	Le taux d'évolution de $7{,}5\,\%$ correspond à un coefficient multiplicateur de :
	\[
		CM = 1+\dfrac{t}{100} = 1+\dfrac{7{,}5}{100} = 1{,}075
	\]

	La suite est donc définie par :
	\[
		u_{\pa{n+1}}=1{,}075 \times u_{\pa{n}} \quad\text{et}\quad u_{\pa{0}}=2\,500
	\]

	$(u_n)$ est donc une \textbf{suite géométrique} de raison $q=1{,}075$.

	\begin{itemize}
		\item $u_{\pa{0}} = 2\,500$
		\item $u_{\pa{1}} = 1{,}075 \times u_{\pa{0}} = 1{,}075 \times 2\,500 = 2687{,}50$
		\item $u_{\pa{2}} = 1{,}075 \times u_{\pa{1}} = 1{,}075 \times 2\,687{,}50 \approx 2\,889{,}06$
		\item $\dots$
	\end{itemize}

	\begin{center}
		\rnc{\arraystretch}{1.4}
		\begin{tabular}{|c|c|c|c|c|c|c|c|}\hline
			\ccg $n$   & $0$      & $1$           & $2$           & $3$           & $\dots$ & \ccg $\textcolor{red}{20}$ & $\dots$ \\ \hline
			\ccg $u_n$ & $2\,500$ & $2\,687{,}50$ & $2\,889{,}06$ & $3\,105{,}74$ & $\dots$ & \ccg $10\,619{,}63$        & $\dots$ \\ \hline
		\end{tabular}
	\end{center}

	Le terme général de la suite s'écrit :
	\[
		u_{\pa{n}}=u_{\pa{0}}\times q^{n} \quad\Lrarr\quad u_{\pa{n}}=2\,500\times 1{,}075^{n}
	\]

	Pour calculer la valeur du capital au bout de $20$ ans ($n = 20$) :
	\[
		u_{\pa{\textcolor{red}{20}}} = 2\,500 \times 1{,}075^{\textcolor{red}{20}} \approx 10\,619{,}63
	\]

	Au bout de $20$ ans, le capital s'élève à environ $10\,619{,}63\eur{}$.

	\begin{center}
		\begin{tikzpicture}[>=Stealth,y=10mm, x=6mm]
			\tikzset{ pt/.style={draw=blue!80!black, fill=blue!60, circle, inner sep=1.3pt} }
			\draw[gray!30, very thin, ystep=0.5, xstep=1] (0, 2) grid (20.5, 11.4);
			\draw[->, thick] (0, 2.0) -- (20.5, 2.0) node[right] {$n$};
			\draw[->, thick] (0, 2.0) -- (0, 11.4) node[above] {$u_{\pa{n}}$ (en $\text{k}\eur$)};
			\foreach \x in {0, 1,..., 20} \draw (\x, 2.025) -- (\x, 1.975) node[below] {$\x$};
			\foreach \y in {2, 2.5,..., 11} \draw (0.1, \y) -- (-0.1, \y) node[left] {$\y$};
			\foreach \n in {0, 1, ..., 20} { \node[pt] at (\n, {2.5 * (1.075^\n)}) {}; }
		\end{tikzpicture}
	\end{center}
\end{ex}

